\documentclass[10pt,a4paper]{amsart}
\usepackage{amsmath,amssymb,amsthm,graphicx}
% Retained minimal preamble for the Stacks Project chapter flat.tex (More on Flatness).
% Verbatim excerpt of the chapter preamble: Xy-pic setup (native preamble lines 14-18) and the native
% theorem declarations with the mathematical macros (native preamble lines 154-217). Copied mechanically;
% nothing added, removed or redefined. The wrapper class and the amsmath/amssymb/amsthm/graphicx base
% packages are supplied by the assembler, so the class/externaldocument part of the native preamble is omitted.
% For commutative diagrams we use Xy-pic
\usepackage[all]{xy}

% We use 2cell for 2-commutative diagrams.
\xyoption{2cell}
% Theorem environments.
%
\theoremstyle{plain}
\newtheorem{theorem}[subsection]{Theorem}
\newtheorem{proposition}[subsection]{Proposition}
\newtheorem{lemma}[subsection]{Lemma}

\theoremstyle{definition}
\newtheorem{definition}[subsection]{Definition}
\newtheorem{example}[subsection]{Example}
\newtheorem{exercise}[subsection]{Exercise}
\newtheorem{situation}[subsection]{Situation}

\theoremstyle{remark}
\newtheorem{remark}[subsection]{Remark}
\newtheorem{remarks}[subsection]{Remarks}

\numberwithin{equation}{subsection}

% Macros
%
\def\lim{\mathop{\mathrm{lim}}\nolimits}
\def\colim{\mathop{\mathrm{colim}}\nolimits}
\def\Spec{\mathop{\mathrm{Spec}}}
\def\Hom{\mathop{\mathrm{Hom}}\nolimits}
\def\Ext{\mathop{\mathrm{Ext}}\nolimits}
\def\SheafHom{\mathop{\mathcal{H}\!\mathit{om}}\nolimits}
\def\SheafExt{\mathop{\mathcal{E}\!\mathit{xt}}\nolimits}
\def\Sch{\mathit{Sch}}
\def\Mor{\mathop{\mathrm{Mor}}\nolimits}
\def\Ob{\mathop{\mathrm{Ob}}\nolimits}
\def\Sh{\mathop{\mathit{Sh}}\nolimits}
\def\NL{\mathop{N\!L}\nolimits}
\def\CH{\mathop{\mathrm{CH}}\nolimits}
\def\proetale{{pro\text{-}\acute{e}tale}}
\def\etale{{\acute{e}tale}}
\def\QCoh{\mathit{QCoh}}
\def\Ker{\mathop{\mathrm{Ker}}}
\def\Im{\mathop{\mathrm{Im}}}
\def\Coker{\mathop{\mathrm{Coker}}}
\def\Coim{\mathop{\mathrm{Coim}}}

% Boxtimes
%
\DeclareMathSymbol{\boxtimes}{\mathbin}{AMSa}{"02}

%
% Macros for moduli stacks/spaces
%
\def\QCohstack{\mathcal{QC}\!\mathit{oh}}
\def\Cohstack{\mathcal{C}\!\mathit{oh}}
\def\Spacesstack{\mathcal{S}\!\mathit{paces}}
\def\Quotfunctor{\mathrm{Quot}}
\def\Hilbfunctor{\mathrm{Hilb}}
\def\Curvesstack{\mathcal{C}\!\mathit{urves}}
\def\Polarizedstack{\mathcal{P}\!\mathit{olarized}}
\def\Complexesstack{\mathcal{C}\!\mathit{omplexes}}
% \Pic is the operator that assigns to X its picard group, usage \Pic(X)
% \Picardstack_{X/B} denotes the Picard stack of X over B
% \Picardfunctor_{X/B} denotes the Picard functor of X over B
\def\Pic{\mathop{\mathrm{Pic}}\nolimits}
\def\Picardstack{\mathcal{P}\!\mathit{ic}}
\def\Picardfunctor{\mathrm{Pic}}
\def\Deformationcategory{\mathcal{D}\!\mathit{ef}}
\usepackage{cleveref}
\crefname{originaldefinition}{original definition}{original definitions}
\Crefname{originaldefinition}{Original definition}{Original definitions}
\makeatletter
\newlabel{categories-definition-initial}{{09WP}{source}{}{sourcepointer-1}{}}
\newlabel{categories-definition-initial@cref}{{[originaldefinition][0][]09WP}{[1][0][]source}{}{}{}}
\makeatother
\crefname{originaldefinition}{original definition}{original definitions}
\Crefname{originaldefinition}{Original definition}{Original definitions}
\makeatletter
\newlabel{categories-definition-codirected}{{04AZ}{source}{}{sourcepointer-2}{}}
\newlabel{categories-definition-codirected@cref}{{[originaldefinition][0][]04AZ}{[1][0][]source}{}{}{}}
\makeatother
\crefname{originallemma}{original lemma}{original lemmas}
\Crefname{originallemma}{Original lemma}{Original lemmas}
\makeatletter
\newlabel{morphisms-lemma-quasi-compact-immersion}{{01RG}{source}{}{sourcepointer-3}{}}
\newlabel{morphisms-lemma-quasi-compact-immersion@cref}{{[originallemma][0][]01RG}{[1][0][]source}{}{}{}}
\makeatother
\crefname{originallemma}{original lemma}{original lemmas}
\Crefname{originallemma}{Original lemma}{Original lemmas}
\makeatletter
\newlabel{morphisms-lemma-scheme-theoretic-image-of-partial-section}{{0CNG}{source}{}{sourcepointer-4}{}}
\newlabel{morphisms-lemma-scheme-theoretic-image-of-partial-section@cref}{{[originallemma][0][]0CNG}{[1][0][]source}{}{}{}}
\makeatother
\title{The category of compactifications of a compactifyable scheme is cofiltered, with the scheme theoretically dense compactifications initial}
\author{The Stacks Project Authors}
\date{stacks-project:flat.tex:local-snapshot-20261011; GNU FDL 1.2 or later}
\begin{document}
\maketitle
\noindent\textit{Editorial note.} All delivered mathematical text is the author's own, copied from the pinned Stacks Project chapter file flat.tex; no statement or proof step was rewritten, reordered or removed.\par\medskip
\noindent\textit{Editorial note.} The preamble is the chapter's own retained minimal preamble (sources/stacks-flat/retained-preamble.tex), a verbatim excerpt of the native Xy-pic setup and of the native theorem declarations and mathematical macros; no macro was invented or redefined, and the wrapper class and the amsmath/amssymb/amsthm/graphicx base packages are supplied by the assembler.\par\medskip
\noindent\textit{Editorial note.} One step of part (b) is flagged by the author as omitted and is kept verbatim in the delivered proof (source lines 9369-9371): "follows from the first (we omit the solution to this categorical exercise)". Part (a) and part (c) are written out in full inside the counted span; the omitted step is the formal consequence of the just-proved cofilteredness recorded in the statement.\par\medskip
\noindent\textit{Editorial note.} Canonical location of the counted proposition: Stacks tag 0ATU (sources/stacks-flat/pinned-original-source.tex line 9316).\par\medskip
\noindent\textit{Allowed standard prerequisites (editorial).}\begin{itemize}
\item Cofiltered (codirected) categories and initial objects and subcategories. Categories, Definitions labelled categories-definition-codirected (Stacks tag 04AZ) and categories-definition-initial (Stacks tag 09WP), invoked at statement line 9325 and proof lines 9334, 9366. Supply the three conditions that cofilteredness amounts to and the two conditions for a full subcategory to be initial; checking those conditions for the category of compactifications is the content of the counted span.
\item Scheme theoretic closure and scheme theoretic image of a partial section. Morphisms, Lemma labelled morphisms-lemma-quasi-compact-immersion (Stacks tag 01RG) and Lemma labelled morphisms-lemma-scheme-theoretic-image-of-partial-section (Stacks tag 0CNG), invoked at proof lines 9364 and 9376. Supply that X is dense and scheme theoretically dense in its scheme theoretic closure, and the topological description of the preimage of j(X) under f; the construction of the dominating compactification and of the equalizer is inside the counted span.
\end{itemize}\medskip
\ifdefined\hypertarget\hypertarget{sourcepointer-1}{}\fi
\noindent\textit{Source reference.} Original definition 09WP: sources/stacks-chow/tags.txt, source lines 12502--12502. Cross-chapter standard prerequisite located in the project tag list staged in the local snapshot; invoked at proof line 9366 of the counted proof. It is retained only as a navigation pointer; every load-bearing step is inside the delivered span.\par\medskip
\ifdefined\hypertarget\hypertarget{sourcepointer-2}{}\fi
\noindent\textit{Source reference.} Original definition 04AZ: sources/stacks-chow/tags.txt, source lines 5634--5634. Cross-chapter standard prerequisite located in the project tag list staged in the local snapshot; invoked at proof line 9334 of the counted proof. It is retained only as a navigation pointer; every load-bearing step is inside the delivered span.\par\medskip
\ifdefined\hypertarget\hypertarget{sourcepointer-3}{}\fi
\noindent\textit{Source reference.} Original lemma 01RG: sources/stacks-chow/tags.txt, source lines 2301--2301. Cross-chapter standard prerequisite located in the project tag list staged in the local snapshot; invoked at proof line 9364 of the counted proof. It is retained only as a navigation pointer; every load-bearing step is inside the delivered span.\par\medskip
\ifdefined\hypertarget\hypertarget{sourcepointer-4}{}\fi
\noindent\textit{Source reference.} Original lemma 0CNG: sources/stacks-chow/tags.txt, source lines 15891--15891. Cross-chapter standard prerequisite located in the project tag list staged in the local snapshot; invoked at proof line 9377 of the counted proof. It is retained only as a navigation pointer; every load-bearing step is inside the delivered span.\par\medskip
\begin{lemma}
\label{lemma-compactifications-cofiltered}
Let $S$ be a quasi-compact and quasi-separated scheme.
Let $X$ be a compactifyable scheme over $S$.
\begin{enumerate}
\item[(a)] The category of compactifications of $X$ over $S$ is
cofiltered.
\item[(b)] The full subcategory consisting of compactifications
$j : X \to \overline{X}$ such that $j(X)$ is dense and
scheme theoretically dense in $\overline{X}$ is initial
(Categories, Definition \ref{categories-definition-initial}).
\item[(c)] If $f : \overline{X}' \to \overline{X}$ is a morphism
of compactifications of $X$ such that $j'(X)$ is dense in $\overline{X}'$,
then $f^{-1}(j(X)) = j'(X)$.
\end{enumerate}
\end{lemma}

\begin{proof}
To prove part (a) we have to check conditions (1), (2), (3) of
Categories, Definition \ref{categories-definition-codirected}.
Condition (1) holds exactly because we assumed that $X$
is compactifyable.
Let $j_i : X \to \overline{X}_i$, $i = 1, 2$ be two compactifications.
Then we can consider the scheme theoretic image $\overline{X}$
of $(j_1, j_2) : X \to \overline{X}_1 \times_S \overline{X}_2$.
This determines a third compactification $j : X \to \overline{X}$
which dominates both $j_i$:
$$
\xymatrix{
(X, \overline{X}_1) & (X, \overline{X}) \ar[l] \ar[r] & (X, \overline{X}_2)
}
$$
Thus (2) holds. Let $f_1, f_2 : \overline{X}_1 \to \overline{X}_2$
be two morphisms between compactifications
$j_i : X \to \overline{X}_i$, $i = 1, 2$.
Let $\overline{X} \subset \overline{X}_1$ be the equalizer of
$f_1$ and $f_2$. As $\overline{X}_2 \to S$ is separated, we see
that $\overline{X}$ is a closed subscheme of $\overline{X}_1$
and hence proper over $S$. Moreover, we obtain an
open immersion $X \to \overline{X}$ because $f_1|_X = f_2|_X = \text{id}_X$.
The morphism $(X \to \overline{X}) \to (j_1 : X \to \overline{X}_1)$
given by the closed immersion $\overline{X} \to \overline{X}_1$
equalizes $f_1$ and $f_2$ which proves condition (3).

\medskip\noindent
Proof of (b). Let $j : X \to \overline{X}$ be a compactification.
If $\overline{X}'$ denotes the scheme theoretic closure of $X$
in $\overline{X}$, then $X$ is dense and scheme theoretically dense
in $\overline{X}'$ by
Morphisms, Lemma \ref{morphisms-lemma-quasi-compact-immersion}.
This proves the first condition of
Categories, Definition \ref{categories-definition-initial}.
Since we have already shown the category of compactifications
of $X$ is cofiltered, the second condition of
Categories, Definition \ref{categories-definition-initial}
follows from the first (we omit the solution to this
categorical exercise).

\medskip\noindent
Proof of (c). After replacing $\overline{X}'$ with the scheme theoretic
closure of $j'(X)$ (which doesn't change the underlying topological space)
this follows from Morphisms, Lemma
\ref{morphisms-lemma-scheme-theoretic-image-of-partial-section}.
\end{proof}

\end{document}
